%%%%%%%%%%%%%%%%%%%%%%%%%%%%%%%%%%%%%%%%%%%%%%%%%%%%%%%%%%%%%%%%%%%%%%
% Lớp 8 | Chương II. Hằng đẳng thức đáng nhớ và ứng dụng | Tổng và hiệu hai lập phương
% Dạng 8D2?8-1: Khai triển, viết dưới dạng tích
%%%%%%%%%%%%%%%%%%%%%%%%%%%%%%%%%%%%%%%%%%%%%%%%%%%%%%%%%%%%%%%%%%%%%%


% nguồn: De_AI_Toan8_cuoi_hoc_ki_i_1 (nhập 21/11/2026)
\begin{bt}%[8D2H8-1]%[ma:b68fdf90]
	Viết các biểu thức sau dưới dạng tích bằng hằng đẳng thức tổng hoặc hiệu hai lập phương:
a) $x^3 + 27$;
b) $8x^3 - 125y^3$.
	\loigiai{
		a) $x^3 + 27 = x^3 + 3^3 = (x + 3)(x^2 - 3x + 9)$. \\
b) $8x^3 - 125y^3 = (2x)^3 - (5y)^3 = (2x - 5y)(4x^2 + 10xy + 25y^2)$.
	}
\end{bt}

% nguồn: De_AI_Toan8_cuoi_hoc_ki_i_1 (nhập 21/11/2026)
\begin{bt}%[8D2H8-1]%[ma:556ace50]
	Viết các biểu thức sau dưới dạng tích:
a) $27x^3 + 1$;
b) $\dfrac{1}{8}x^3 - 64y^3$.
	\loigiai{
		a) $27x^3 + 1 = (3x)^3 + 1^3 = (3x + 1)(9x^2 - 3x + 1)$. \\
b) $\dfrac{1}{8}x^3 - 64y^3 = \left(\dfrac{1}{2}x\right)^3 - (4y)^3 = \left(\dfrac{1}{2}x - 4y\right)\left(\dfrac{1}{4}x^2 + 2xy + 16y^2\right)$.
	}
\end{bt}

% nguồn: De_AI_Toan8_cuoi_hoc_ki_i_1 (nhập 21/11/2026)
